\section{Conclusion}
\label{sec:conclusion}

The inferential target decides what additional cases buy: more cases shrink the counting error of the current value and of the realised mean of the observed periods, but they do not remove the between-period variation that limits estimation of the process mean. When reporting fluctuates from week to week, that fluctuation leaves an error floor on the current value, and single-period second moments cannot separate it from transmission variation. Cross-sectional dispersion diagnostics therefore measure a count-independent relative dispersion and the counting noise around it rather than a precision limit or a threshold for allocating surveillance effort, and separating under-reporting, superspreading and structured transmission requires temporal models or external data on reporting.
